% STATUS: COMPLETE DRAFT
% Chapter 4 : exercises (solutions in Appendix A). Every number is checked in scripts/verify.py
\bichapter{Exercises}{Esercizi}\label{ch:exercises}

\begin{bi}
\colheads
Try each exercise before looking at the solutions in Appendix~\ref{app:solutions}. For the proofs, follow the three steps of section~\ref{sec:howto}: figure, hypothesis and thesis, then the chain of reasons.
\IT
Prova ogni esercizio prima di guardare le soluzioni nell'appendice~\ref{app:solutions}. Per le dimostrazioni segui i tre passi del paragrafo~\ref{sec:howto}: figura, ipotesi e tesi, poi la catena dei motivi.
\end{bi}

\begin{exercise}{A, angles}{A, angoli}
\begin{enumerate}[label=\textbf{A\arabic*.}, leftmargin=2.2em]
\item Find the complement, the supplement and the explement of $72^\circ$.
\item What angle do the hands of a clock make at 7:20?
\item Of 40 students, 14 come to school by bus, 10 by bike, 8 by car and 8 on foot. Find the angles of the pie chart.
\item (a) Which angle is equal to its complement? (b) Which angle is three times its supplement?
\item Two vertical angles measure $4x + 10^\circ$ and $2x + 50^\circ$. Find $x$ and all four angles of the X.
\end{enumerate}
\tcblower
\begin{enumerate}[label=\textbf{A\arabic*.}, leftmargin=2.2em]
\item Trova il complementare, il supplementare e l'esplementare di $72^\circ$.
\item Che angolo formano le lancette dell'orologio alle 7:20?
\item Su 40 studenti, 14 vengono a scuola in autobus, 10 in bici, 8 in auto e 8 a piedi. Trova gli angoli dell'areogramma.
\item (a) Quale angolo è congruente al suo complementare? (b) Quale angolo è il triplo del suo supplementare?
\item Due angoli opposti al vertice misurano $4x + 10^\circ$ e $2x + 50^\circ$. Trova $x$ e tutti e quattro gli angoli della X.
\end{enumerate}
\end{exercise}

\begin{exercise}{B, parallel lines}{B, rette parallele}
\begin{enumerate}[label=\textbf{B\arabic*.}, leftmargin=2.2em]
\item In Figure~\ref{fig:exfig}a, $r \parallel s$ and one angle is $118^\circ$. Find the other seven and, where the pair has a name, say what pair each forms with the $118^\circ$ angle.
\item Two corresponding angles measure $5x - 20^\circ$ and $3x + 30^\circ$ ($r \parallel s$). Find $x$ and the angles.
\item True or false? (a) Alternate exterior angles add up to $180^\circ$. (b) If $r \parallel s$, same side interior angles add up to $180^\circ$. (c) If the alternate interior angles are equal, the lines are parallel.
\end{enumerate}
\tcblower
\begin{enumerate}[label=\textbf{B\arabic*.}, leftmargin=2.2em]
\item Nella figura~\ref{fig:exfig}a, $r \parallel s$ e un angolo misura $118^\circ$. Trova gli altri sette e, quando la coppia ha un nome, di' che coppia forma ciascuno con l'angolo di $118^\circ$.
\item Due angoli corrispondenti misurano $5x - 20^\circ$ e $3x + 30^\circ$ ($r \parallel s$). Trova $x$ e gli angoli.
\item Vero o falso? (a) Gli angoli alterni esterni hanno somma $180^\circ$. (b) Se $r \parallel s$, gli angoli coniugati interni hanno somma $180^\circ$. (c) Se gli angoli alterni interni sono congruenti, le rette sono parallele.
\end{enumerate}
\end{exercise}

\begin{exercise}{C, triangles and altitudes}{C, triangoli e altezze}
\begin{enumerate}[label=\textbf{C\arabic*.}, leftmargin=2.2em]
\item The angles of a triangle are in the ratio $2 : 3 : 4$. Find them.
\item In an isosceles triangle the vertex angle is twice each base angle. Find the three angles. What kind of triangle is it?
\item Can a triangle have two obtuse angles? Two right angles? Explain.
\item A triangle has a side of 10\,cm with altitude 6\,cm. Another side is 12\,cm. How long is the altitude to that side?
\item Draw an obtuse triangle and its three altitudes. Where is the orthocentre?
\end{enumerate}
\tcblower
\begin{enumerate}[label=\textbf{C\arabic*.}, leftmargin=2.2em]
\item Gli angoli di un triangolo stanno nel rapporto $2 : 3 : 4$. Trovali.
\item In un triangolo isoscele l'angolo al vertice è il doppio di ciascun angolo alla base. Trova i tre angoli. Che triangolo è?
\item Un triangolo può avere due angoli ottusi? Due angoli retti? Spiega.
\item Un triangolo ha un lato di 10\,cm con altezza 6\,cm. Un altro lato misura 12\,cm. Quanto è lunga l'altezza relativa a quel lato?
\item Disegna un triangolo ottusangolo e le sue tre altezze. Dove si trova l'ortocentro?
\end{enumerate}
\end{exercise}

\begin{exercise}{D, congruence proofs}{D, dimostrazioni di congruenza}
\begin{enumerate}[label=\textbf{D\arabic*.}, leftmargin=2.2em]
\item $ABC$ is isosceles with base $BC$. Take $D$ on $AB$ and $E$ on $AC$ with $\seg{AD} \cong \seg{AE}$. Prove $\seg{BE} \cong \seg{CD}$.
\item $ABC$ is equilateral. Take $D$ on $AB$, $E$ on $BC$, $F$ on $CA$ with $\seg{AD} \cong \seg{BE} \cong \seg{CF}$ (Figure~\ref{fig:exfig}b). Prove that $DEF$ is equilateral.
\item $ABC$ is isosceles with base $BC$ and $M$ is the midpoint of $BC$. Prove that $AM \perp BC$, using the third criterion.
\item $M$ is the midpoint of $AB$ and $r$ is a line through $M$. $H$ and $K$ are the feet of the perpendiculars from $A$ and $B$ to $r$ (Figure~\ref{fig:exfig}c). Prove $\seg{AH} \cong \seg{BK}$.
\item Find the mistake: ``$\av{A} \cong \av{D}$, $\av{B} \cong \av{E}$, $\av{C} \cong \av{F}$, so $\tri{ABC} \cong \tri{DEF}$.''
\end{enumerate}
\tcblower
\begin{enumerate}[label=\textbf{D\arabic*.}, leftmargin=2.2em]
\item $ABC$ è isoscele di base $BC$. Prendi $D$ su $AB$ ed $E$ su $AC$ con $\seg{AD} \cong \seg{AE}$. Dimostra che $\seg{BE} \cong \seg{CD}$.
\item $ABC$ è equilatero. Prendi $D$ su $AB$, $E$ su $BC$, $F$ su $CA$ con $\seg{AD} \cong \seg{BE} \cong \seg{CF}$ (figura~\ref{fig:exfig}b). Dimostra che $DEF$ è equilatero.
\item $ABC$ è isoscele di base $BC$ e $M$ è il punto medio di $BC$. Dimostra che $AM \perp BC$, usando il terzo criterio.
\item $M$ è il punto medio di $AB$ e $r$ è una retta per $M$. $H$ e $K$ sono i piedi delle perpendicolari da $A$ e da $B$ a $r$ (figura~\ref{fig:exfig}c). Dimostra che $\seg{AH} \cong \seg{BK}$.
\item Trova l'errore: «$\av{A} \cong \av{D}$, $\av{B} \cong \av{E}$, $\av{C} \cong \av{F}$, quindi $\tri{ABC} \cong \tri{DEF}$.»
\end{enumerate}
\end{exercise}

\begin{figure}[htbp]
\centering
% === PRE-COMPUTATION ===
% (a) r: y = 1.8, s: y = 0, transversal through Q = (0,0) at 118 deg (so the angle above s on the right is 118):
%     P = (1.8/tan118, 1.8) = (-0.957, 1.8).  VERIFIED: 180 - 118 = 62
% (b) equilateral A(0,0) B(3.2,0) C(1.6,2.771); t = 0.3: D(0.96,0) E(2.72,0.831) F(1.12,1.940)
%     VERIFIED (scaled model in verify.py, side 4): DE = EF = FD
% (c) A(0,0) B(4,0) M(2,0); r through M at 35 deg; H(0.658,-0.940) K(3.342,0.940). VERIFIED: AH = BK = 1.1472
\begin{tikzpicture}[scale=0.95]
  \begin{scope}
    \coordinate (Q) at (0,0); \coordinate (P) at (-0.957,1.8);
    \draw[side] (-2.2,1.8) -- (1.6,1.8) node[right] {$r$};
    \draw[side] (-2.2,0) -- (1.6,0) node[right] {$s$};
    \draw[side] ($(Q)!-0.45!(P)$) -- ($(Q)!1.45!(P)$) node[above] {$t$};
    \fill[angA!35] (Q) -- ++(0:0.45) arc (0:118:0.45) -- cycle;
    \node[font=\small, angA] at ($(Q)+(59:0.75)$) {$118^\circ$};
    \node[font=\small\bfseries] at (-0.3,-1.25) {(a)};
  \end{scope}
  \begin{scope}[shift={(4.3,-0.2)}]
    \coordinate (A) at (0,0); \coordinate (B) at (3.2,0); \coordinate (C) at (1.6,2.771);
    \coordinate (D) at (0.96,0); \coordinate (E) at (2.72,0.831); \coordinate (F) at (1.12,1.940);
    \draw[side] (A) -- (B) -- (C) -- cycle;
    \fill[angB!15] (D) -- (E) -- (F) -- cycle; \draw[side, angB] (D) -- (E) -- (F) -- cycle;
    \draw[tick1] (A) -- (D); \draw[tick1] (B) -- (E); \draw[tick1] (C) -- (F);
    \node[below left] at (A) {$A$}; \node[below right] at (B) {$B$}; \node[above] at (C) {$C$};
    \node[below] at (D) {$D$}; \node[right] at (E) {$E$}; \node[left] at (F) {$F$};
    \node[font=\small\bfseries] at (1.6,-1.05) {(b)};
  \end{scope}
  \begin{scope}[shift={(9.2,0.4)}]
    \coordinate (A) at (0,0); \coordinate (B) at (4,0); \coordinate (M) at (2,0);
    \coordinate (H) at (0.658,-0.940); \coordinate (K) at (3.342,0.940);
    \draw[side] ($(M)+(215:2.6)$) -- ($(M)+(35:2.6)$) node[above right] {$r$};
    \draw[side] (A) -- (B);
    \draw[side, angB] (A) -- (H); \draw[side, angB] (B) -- (K);
    \rang{A}{H}{M} \rang{B}{K}{M}
    \node[left] at (A) {$A$}; \node[right] at (B) {$B$}; \node[below] at ($(M)+(0.1,-0.05)$) {$M$};
    \node[below] at (H) {$H$}; \node[above] at (K) {$K$};
    \node[font=\small\bfseries] at (2,-1.65) {(c)};
  \end{scope}
\end{tikzpicture}
\bicap{Figures for exercises B1, D2 and D4.}{Figure per gli esercizi B1, D2 e D4.}
\label{fig:exfig}
\end{figure}
